The fundamental-mode RF system of an electron storage ring and the stored beam form one coupled dynamical system whose behavior changes with beam current and with the settings of the RF feedback. We ask whether closed-loop measurements made during operation can identify a compact model that predicts these changes. From eighteen closed-loop transfer-function records of the SPEAR3 RF station, taken with its existing noise-injection routine at beam currents from 20 to \SI{500}{mA} and at several loop gains and phases, we identify a physics-informed two-quadrature model of beam, cavities, klystron and feedback. Single records, measured with a one-dimensional excitation, do not determine the klystron's response asymmetry. Records at several loop settings, together with the band of the coupled beam mode, which the routine itself discards, determine its compression but leave its AM/PM conversion correlated with the loop phase; records frozen at trips of the same klystron supply it ($k_\mathrm{AP}=\tripKAP\pm\tripKAPErr$). Of the competing explanations for the large change of the response with beam current---effective linear parameters, unequal cavity detunings, beam dynamics at fixed hardware, and a compressed klystron---only the last describes all records with one set of station constants---at the noise level in the band of the coupled beam mode and to within a few percent across the loop band, with \nparSel{} parameters for \nPoints{} data points---and these constants agree with independent station data. The klystron's compression, its amplitude gain relative to its phase gain, falls from \gaSelLow{} at \SI{22}{mA} to \gaSelHighLo{} at \SI{500}{mA} as the tube, held at a fixed drive, works closer to saturation; at full output the trip records give \tripGAHigh--\tripGAFull. The model predicts records at beam currents it was not trained on almost at in-sample quality, where interpolated effective parameters fail. Its out-of-band prediction of the coupled mode isolates the synchrotron focusing as the one beam parameter the loop band cannot determine; identified from the band of the coupled mode, it implies an incoherent synchrotron frequency \fsRatio\% above the design value, a prediction that a low-current measurement can test.
